%---------------------------------------------------------------------------
\chapter{連続システムの離散化}
%---------------------------------------------------------------------------
\section{状態方程式の離散化}
%---------------------------------------------------------------------------
システム
\begin{align}
 \frac{dx}{dt} &= A_c x(t) + B_c u(t) \label{c_state_eq} \\
 y(t)          &= C x(t) + D u(t) \label{c_output_eq}
\end{align}
を離散化することを考える．いま
サンプリングインターバルを$T$とし，時刻$kT$における状態を$x_k$, 出力を
$y_k$，入力$u(t)$は時刻$kT \leq t < (k+1)T$において$u(t) = u_k$とする．
(Fig.\ref{fig_discrete_signal}参照)

\begin{figure}[htbp]
 \begin{center}
  \input{figure/discretize.pstex_t} % scale = 100%
  \caption{システムの離散化}
  \label{fig_discrete_signal}
 \end{center}
\end{figure}
\newpage

このようなシステムの$T$時間後（１サンプル時間後）の応答は，初期値$x(kT)$,積分区間$T$として，入力のある状態方程式の応答(\ref{response_w_input})を用いれば
\begin{equation}
 x(kT+T) = e^{A_c T} x(kT) 
  + \int_0^T e^{A_c(T - \tau)}B_c u (kT + \tau) d \tau
\end{equation}
と表される．ここで仮定より$u(kT+\tau)$, $(0\leq\tau<T)$は一定値$u_k$であるから，離散系の状態方程式は
\begin{align}
 x_{k+1} &= e^{A_cT} x_k 
  + \left( \int_0^T e^{A_c(T-\tau)} B_c d \tau \right) u_k \nonumber \\
 &= A x_k + B u_k
\end{align}
と表される．また出力方程式は明かに
\begin{equation}
 y_k = C x_k + D u_k 
\end{equation}
となる．これが離散化されたシステムの状態方程式と出力方程式である．

離散化されたシステムの$A$行列と元の連続時間システムの$A_c$行列の固有値に
は次のような関係がある．いま$A_c$の固有値を$\lambda_{c1}, \lambda_{c2},
\cdots \lambda_{cn}$とし，正則な複素行列$S$を
\begin{align}
 S^{-1} A_c S 
  &=
  \left(\begin{array}{cccc}
         \lambda_{c1} & 0            & \cdots & 0 \\
	 0            & \lambda_{c2} &        & 0 \\
	 \vdots       &              & \ddots & \\
	 0            & \cdots       &        & \lambda_{cn} 
	\end{array}\right) \nonumber \\
 &= \Lambda
\end{align}
を満たすように選ぶことができると仮定する(ジョルダン形式になる場合も後の
議論はほとんど同じになる)．この$\Lambda$を用いれば
\begin{align}
 A &=  e^{A_cT} \nonumber \\
 &= S e^{\Lambda T} S^{-1} \nonumber \\
 &=
  S \left(\begin{array}{cccc}
           e^{\lambda_{c1} T} & 0                  & \cdots & 0 \\
	   0                  & e^{\lambda_{c2} T} &        & 0 \\
	   \vdots             &                    & \ddots & \\
	   0                  & \cdots             & & e^{\lambda_{cn} T}
	  \end{array}\right)
  S^{-1}
\end{align}
となり，離散化されたシステムの$A$行列の固有値は$e^{\lambda_{ci} T}$とな
る．よって連続時間システムの$A_c$行列の固有値の実部が負であるならば，離散
化されたシステムの$A$行列の固有値の絶対値は1未満となる．これは後で述べる
ように連続システムが安定であれば離散化されたシステムも安定となることを示
している．

%---------------------------------------------------------------------------
\section{連続時間コントローラの離散化}
%---------------------------------------------------------------------------
連続系に対して連続コントローラを設計した場合，これをディジタル計算機上で
実現しようとするときにはこれを離散化しなければならない．もし，連続コント
ローラが状態方程式で表されている(例えば状態観測器と状態フィードバック)
ならばこれを離散化すれば良い．

特殊なシステムの場合は，先の離散化以外の方法で制御できる場合がある．
例えば機械系の制御の場合はつぎのようにオブザーバを用いないで制御することができる．
一般に機械系の状態方程式は連続系において
\begin{align}
 \frac{d}{dt}
  \left(\begin{array}{c}
         x_1 \\
	 x_2 
	\end{array}\right)
  &=
  \left(\begin{array}{cc}
         0   & I \\
	 K_1 & K_2
	\end{array}\right)
  \left(\begin{array}{c}
         x_1 \\
	 x_2 
	\end{array}\right)
  +
  \left(\begin{array}{c}
         0 \\
	 K_3
	\end{array}\right)
  u \\
 y &=
  \left(\begin{array}{cc}
         I & 0 
	\end{array}\right)
  \left(\begin{array}{c}
         x_1 \\
	 x_2 
	\end{array}\right)
\end{align}
と表されることが多い．ここで
\begin{align}
 x_1 &= y \\
 x_2 &= \frac{dy}{dt}
\end{align}
であることから，状態フィードバックは
\begin{equation}
 u(t) = F_1 y + F_2 \frac{dy}{dt}
\end{equation}
で表される．ここで離散系の制御則として次の状態フィードバックを用いること
を考える．
\begin{equation}
 u_k = F_1 y(kT) + F_2 \left.\frac{dy}{dt} \right\vert_{t = kT}
\end{equation}
このときは$\frac{dy}{dt}$を
\begin{equation}
 \left.\frac{dy}{dt}\right\vert_{t = kT} = \frac{y_k - y_{k-1}}{T}
\end{equation}
で近似すれば状態観測器を用いることなく状態フィードバックをすることが可能
になる．出力の観測に雑音が多い場合には連続系の近似微分器$G_d(s)$
\begin{equation}
 G_d(s) = \frac{s}{Ts + 1}
\end{equation}
を離散化したものを用いて出力を近似微分した方が良い．

いままで述べた連続時間コントローラの離散化は，あくまでサンプリングインター
バル$T$が十分小さいとして連続時間コントローラをディジタル計算機上で実現
するものであった．しかし，サンプリングインターバルが長くなったりすると，連
続時間コントローラと同等の性能が得られない場合がある．このような場合には，
何らかの意味でディジタルコントローラの性能が連続時間コントローラの性
能に近付くように，コントローラの離散化を工夫する必要がある．この問題はディ
ジタル再設計の問題として議論されている．しかし，サンプリングインターバル
が閉ループ系の応答に対して十分短ければ，無理に再設計をする必要はないようで
ある．ディジタル再設計については成書を参考にされたい．
